\Needspace{8\baselineskip}
\SourceLeaf{320}
\Needspace{8\baselineskip}
\Article{277}
The rules for finding all the forms in I and II of the preceding article follow directly from their description, so a few examples will be enough. For \(D=1\), the choices of signs give these six forms I.
\[
\left(\begin{smallmatrix}0,&1,&0\\0,&\pm1,&0\end{smallmatrix}\right),\quad
\left(\begin{smallmatrix}0,&1,&\pm1\\0,&\pm1,&0\end{smallmatrix}\right)
\]
In forms II, \(a\) and \(A''\) can have no values other than \(+1\) and \(-1\). For each of the four resulting combinations, \(b''\), \(B\) and \(B'\) must be set \(=0\), giving four forms
\[
\left(\begin{smallmatrix}1,&-1,&1\\0,&0,&0\end{smallmatrix}\right),\quad
\left(\begin{smallmatrix}-1,&1,&1\\0,&0,&0\end{smallmatrix}\right),\quad
\left(\begin{smallmatrix}1,&1,&-1\\0,&0,&0\end{smallmatrix}\right),\quad
\left(\begin{smallmatrix}-1,&-1,&-1\\0,&0,&0\end{smallmatrix}\right)
\]
Similarly, for \(D=-1\) we have six forms I and four forms II,
\[
\begin{gathered}
\left(\begin{smallmatrix}0,&-1,&0\\0,&\pm1,&0\end{smallmatrix}\right),\quad
\left(\begin{smallmatrix}0,&-1,&\pm1\\0,&\pm1,&0\end{smallmatrix}\right);\quad
\left(\begin{smallmatrix}1,&-1,&-1\\0,&0,&0\end{smallmatrix}\right),\\
\left(\begin{smallmatrix}-1,&1,&-1\\0,&0,&0\end{smallmatrix}\right),\quad
\left(\begin{smallmatrix}-1,&-1,&1\\0,&0,&0\end{smallmatrix}\right),\quad
\left(\begin{smallmatrix}1,&1,&1\\0,&0,&0\end{smallmatrix}\right)
\end{gathered}
\]
For \(D=2\) we obtain six forms I
\[
\left(\begin{smallmatrix}0,&2,&0\\0,&\pm1,&0\end{smallmatrix}\right),\quad
\left(\begin{smallmatrix}0,&2,&\pm1\\0,&\pm1,&0\end{smallmatrix}\right)
\]
and eight forms II
\[
\begin{gathered}
\left(\begin{smallmatrix}1,&-1,&2\\0,&0,&0\end{smallmatrix}\right),\quad
\left(\begin{smallmatrix}-1,&1,&2\\0,&0,&0\end{smallmatrix}\right),\quad
\left(\begin{smallmatrix}1,&1,&-2\\0,&0,&0\end{smallmatrix}\right),\quad
\left(\begin{smallmatrix}-1,&-1,&-2\\0,&0,&0\end{smallmatrix}\right)\\
\left(\begin{smallmatrix}1,&-2,&1\\0,&0,&0\end{smallmatrix}\right),\quad
\left(\begin{smallmatrix}-1,&2,&1\\0,&0,&0\end{smallmatrix}\right),\quad
\left(\begin{smallmatrix}1,&2,&-1\\0,&0,&0\end{smallmatrix}\right),\quad
\left(\begin{smallmatrix}-1,&-2,&-1\\0,&0,&0\end{smallmatrix}\right)
\end{gathered}
\]
The number of classes arising from these forms in these three cases is, however, much smaller than the number of forms. For it is easily confirmed:

I. The form \(\left(\begin{smallmatrix}0,&1,&0\\0,&1,&0\end{smallmatrix}\right)\) passes into
\[
\left(\begin{smallmatrix}0,&1,&0\\0,&-1,&0\end{smallmatrix}\right),\quad
\left(\begin{smallmatrix}0,&1,&1\\0,&\pm1,&0\end{smallmatrix}\right),\quad
\left(\begin{smallmatrix}0,&1,&-1\\0,&\pm1,&0\end{smallmatrix}\right),\quad
\left(\begin{smallmatrix}1,&1,&-1\\0,&0,&0\end{smallmatrix}\right)
\]
by the respective substitutions
\[
\begin{array}{rrr|rrr|rrr|rrr}
1,&0,&0&0,&0,&1&0,&0,&1&1,&0,&-1\\
0,&1,&0&0,&1,&-1&0,&1,&1&1,&1,&-1\\
0,&0,&-1&\pm1,&1,&0&\pm1,&-1,&-1&0,&-1,&1
\end{array}
\]
while the form \(\left(\begin{smallmatrix}1,&1,&-1\\0,&0,&0\end{smallmatrix}\right)\) passes into \(\left(\begin{smallmatrix}1,&-1,&1\\0,&0,&0\end{smallmatrix}\right)\), \(\left(\begin{smallmatrix}-1,&1,&1\\0,&0,&0\end{smallmatrix}\right)\) merely by permuting the indeterminates. Thus those ten ternary forms of determinant \(1\) reduce to these two
\SourceLeaf{321}
\(\left(\begin{smallmatrix}0,&1,&0\\0,&1,&0\end{smallmatrix}\right)\), \(\left(\begin{smallmatrix}-1,&-1,&-1\\0,&0,&0\end{smallmatrix}\right)\); instead of the first, if preferred, we can also take \(\left(\begin{smallmatrix}1,&0,&0\\1,&0,&0\end{smallmatrix}\right)\). Since the former form is indefinite and the latter definite, it is clear that every indefinite ternary form of determinant \(1\) is equivalent to the form \(xx+2yz\), and every definite one to \(-xx-yy-zz\).

II. In exactly the same way we find that every indefinite ternary form of determinant \(-1\) is equivalent to the form \(-xx+2yz\), and every definite one to \(xx+yy+zz\).

III. For determinant \(2\), the second, sixth and seventh of the eight forms (II) can be discarded immediately, since they arise from the first merely by permuting the indeterminates. For the same reason we can also discard the fifth, which arises in this way from the third, and the eighth, which arises from the fourth. The three remaining forms, together with the six forms I, constitute three classes. For \(\left(\begin{smallmatrix}0,&2,&0\\0,&1,&0\end{smallmatrix}\right)\) passes into \(\left(\begin{smallmatrix}0,&2,&0\\0,&-1,&0\end{smallmatrix}\right)\) by the substitution \(\left\{\begin{smallmatrix}1,&0,&0\\0,&1,&0\\0,&0,&-1\end{smallmatrix}\right\}\), and the form \(\left(\begin{smallmatrix}1,&1,&-2\\0,&0,&0\end{smallmatrix}\right)\) into
\[
\begin{gathered}
\left(\begin{smallmatrix}0,&2,&1\\0,&1,&0\end{smallmatrix}\right),\quad
\left(\begin{smallmatrix}0,&2,&1\\0,&-1,&0\end{smallmatrix}\right),\quad
\left(\begin{smallmatrix}0,&2,&-1\\0,&1,&0\end{smallmatrix}\right),\\
\left(\begin{smallmatrix}0,&2,&-1\\0,&-1,&0\end{smallmatrix}\right),\quad
\left(\begin{smallmatrix}1,&-1,&2\\0,&0,&0\end{smallmatrix}\right)
\end{gathered}
\]
by the respective substitutions
\[
\begin{array}{rrr|rrr|rrr|rrr|rrr}
1,&0,&1&1,&0,&-1&1,&0,&0&1,&0,&0&1,&0,&0\\
1,&2,&0&1,&2,&0&1,&2,&-1&1,&2,&1&0,&1,&2\\
1,&1,&0&1,&1,&0&1,&1,&-1&1,&1,&1&0,&1,&1
\end{array}
\]
Thus every ternary form of determinant \(2\) can be reduced to one of these three
\[
\left(\begin{smallmatrix}0,&2,&0\\0,&-1,&0\end{smallmatrix}\right),\quad
\left(\begin{smallmatrix}1,&1,&-2\\0,&0,&0\end{smallmatrix}\right),\quad
\left(\begin{smallmatrix}-1,&-1,&-2\\0,&0,&0\end{smallmatrix}\right)
\]
Instead of the first, if preferred, we can also take \(\left(\begin{smallmatrix}2,&0,&0\\1,&0,&0\end{smallmatrix}\right)\). Clearly every definite ternary form must be equivalent to the third, \(-xx-yy-2zz\), since the first two are indefinite. Every indefinite form is equivalent to the first or second: to the first, \(2xx+2yz\), if its first, second and third coefficients are all even (since it is easy to see that any substitution takes such a form into a form of the same kind, so it cannot be equivalent to the second form); but to the second, \(xx+yy-2zz\), if its first, second and third coefficients are not all even, but one, two or all are odd (for the same reason, the first form \(2xx+2yz\) cannot be transformed into such a form by any substitution).
\SourceLeaf{322}
So what happened in the examples of Articles 273, 274, namely that the definite form
$\left(\begin{smallmatrix}19,21,50\\15,28,1\end{smallmatrix}\right)$
of determinant $-1$ reduced to $xx+yy+zz$, and the indefinite form
$\left(\begin{smallmatrix}10,26,2\\7,0,4\end{smallmatrix}\right)$
of determinant $2$ to $2xx-2yz$, or, which amounts to the same thing, to $2xx+2yz$, could have been predicted in advance from the preceding investigations.

\Needspace{8\baselineskip}
\Article{278}
A ternary form with indeterminates $x,x',x''$ \textit{represents} both numbers, by assigning specific values to $x,x',x''$, and binary forms, by substitutions of the kind
\[
x=mt+nu,\qquad x'=m't+n'u,\qquad x''=m''t+n''u
\]
where $m,n,m'$, etc.\ denote fixed numbers, and $t,u$ the indeterminates of the represented form. A complete theory of ternary forms would therefore require solutions to the following problems: I.\ Find all representations of a given number by a given ternary form. II.\ Find all representations of a given binary form by a given ternary form. III.\ Decide whether two given ternary forms of the same determinant are equivalent, and if they are, find all transformations of one into the other. IV.\ Decide whether a given ternary form contains another given form of larger determinant, and if it does, find all transformations of the former into the latter. We will discuss these problems, which are much more difficult than the analogous problems for binary forms, more fully elsewhere. Here we restrict our investigation to showing how the first problem can be reduced to the second, and the second to the third. We will show how to solve the third in some very simple cases that shed particular light on the theory of binary forms. We exclude the fourth entirely here.

\Needspace{8\baselineskip}
\Article{279}
\textsc{Lemma.} \textit{Given any three integers $a,a',a''$ (not all $=0$): find six others $B,B',B'',C,C',C''$ such that}
\[
B'C''-B''C'=a,\qquad B''C-BC''=a',\qquad BC'-B'C=a''
\]

\textit{Solution.} Let $\alpha$ be the greatest common divisor of $a,a',a''$. and choose integers $A,A',A''$ so that
\[
Aa+A'a'+A''a''=\alpha
\]
\SourceLeaf{323}
Next choose three integers $\mathfrak{C},\allowbreak \mathfrak{C}',\allowbreak \mathfrak{C}''$ arbitrarily, with the sole condition that the three numbers
$\mathfrak{C}'A''-\mathfrak{C}''A',\allowbreak \ \mathfrak{C}''A-\mathfrak{C}A'',\allowbreak \
\mathfrak{C}A'-\mathfrak{C}'A$, which we will denote, respectively, by $b,b',b''$, and their greatest common divisor by $\beta$, are not all $=0$. Then set
\[
a'b''-a''b'=\alpha\beta C,\qquad
a''b-ab''=\alpha\beta C',\qquad
ab'-a'b=\alpha\beta C''
\]
Clearly $C,C',C''$ will be integers. Finally, choose integers
$\mathfrak{B},\allowbreak \mathfrak{B}',\allowbreak \mathfrak{B}''$ so that
\[
\mathfrak{B}b+\mathfrak{B}'b'+\mathfrak{B}''b''=\beta
\]
put
\[
\mathfrak{B}a+\mathfrak{B}'a'+\mathfrak{B}''a''=h
\]
and set
\[
B=\alpha\mathfrak{B}-hA,\qquad
B'=\alpha\mathfrak{B}'-hA',\qquad
B''=\alpha\mathfrak{B}''-hA''
\]
These values of $B,B',B'',C,C',C''$ will satisfy the required equations.

For we find
\[
\begin{gathered}
aB+a'B'+a''B''=0\\
bA+b'A'+b''A''=0\quad\text{so}\quad bB+b'B'+b''B''=\alpha\beta
\end{gathered}
\]
Now the values of $C',C''$ give
\[
\begin{aligned}
\alpha\beta(B'C''-B''C')
&=ab'B'-a'bB'-a''bB''+ab''B''\\
&=a(bB+b'B'+b''B'')-b(aB+a'B'+a''B'')=\alpha\beta a
\end{aligned}
\]
and hence $B'C''-B''C'=a$; similarly we find $B''C-BC''=a'$, $BC'-B'C=a''$. \textit{Q.E.F.} --- We must omit here the analysis by which this solution was found, as well as the method for finding all solutions from one.
